% TEMPLATE-MILC-v3.tex - plantilla maestra UNICA de las guias MILC v3.
%
% La usan las cuatro familias de guias, cada una con su driver:
%   · Grados 6-11          -> scripts/build-guias-g11.py
%   · Bebras               -> scripts/build-guias-bebras.py
%   · Semillero            -> scripts/build-guias-semillero.py
%   · Territorio interior  -> scripts/build-guias-territorio-interior.py
%
% Antes habia tres copias casi identicas (~400 lineas iguales, 6 distintas):
% cada mejora habia que aplicarla tres veces, y en la practica no se hacia
% (el motor de recursos vivia solo en dos de ellas). Lo que varia entre
% familias esta parametrizado en 6 placeholders que cada driver llena
% (se nombran aqui SIN su sintaxis real: el builder reemplaza por texto plano,
% tambien dentro de los comentarios, y un valor multilinea rompe el preambulo):
%
%   HEADER_IZQ        encabezado izquierdo de pagina
%   HEADER_DER        encabezado derecho de pagina
%   PORTADA_ETIQUETA  rotulo sobre el numero grande (GRADO/MOMENTO/MODULO)
%   PORTADA_SERIE     linea de serie de la portada
%   PORTADA_ID        identificador de la guia en la portada
%   PORTADA_CIRCULO   el circulito de grado (°), o vacio si no aplica
%   PDF_TITULO        titulo en texto plano (sin \textbf ni comillas TeX) para
%                     los metadatos del PDF (/Title); lo produce lib_xelatex.texto_pdf
%
% Composicion (auditoria 2026-09-03): las bandas de estacion piden espacio
% (\Needspace*) y prohiben el salto justo despues (\nopagebreak), asi no quedan
% huerfanas al pie; las cajas largas (softbox, drawbox, infoband, puentebox) son
% `breakable`, asi una caja de media pagina ya no arrastra todo a la siguiente
% dejando la anterior al 40 %. Todo texto sobre mostaza o verde va en milcNegro
% (blanco sobre #E5B400 da 1,93:1 y sobre #58A923 2,95:1; ninguno pasa WCAG AA).
% El resto de claves las documenta content/guias/_SCHEMA.md. El script valida
% que no queden placeholders sin reemplazar antes de escribir el archivo final.

% Etiquetado PDF (latex-lab): /Lang, estructura y texto alternativo de las
% figuras, para lectores de pantalla. Debe ir ANTES de \documentclass.
\DocumentMetadata{lang=es-CO,tagging=on}
\documentclass[11pt,letterpaper]{article}
% headheight=24pt: la cabecera derecha lleva dos lineas (identidad de la guia
% y estacion actual). headsep se acorta en la misma medida para que el bloque
% de texto quede exactamente donde estaba con headheight=12pt.
\usepackage[margin=2.0cm,headheight=24pt,headsep=13pt]{geometry}
\usepackage[table]{xcolor}
\usepackage{fontspec}
\usepackage[spanish]{babel}
\usepackage{tikz}
% Librerías TikZ para los diagramas de las guías (bloque `recursos:`):
%  · babel  → babel en español vuelve activos caracteres como «>», y TikZ los usa
%             en sus opciones (>=stealth, ->). Sin esta librería, cualquier
%             diagrama con flechas revienta con «\language@active@arg> has an
%             extra }». Debe ir ANTES de que se dibuje nada.
%  · positioning        → sintaxis «below left=2pt and 3pt».
%  · decorations.pathreplacing → llaves ({brace}) para señalar rangos.
\usetikzlibrary{babel,positioning,decorations.pathreplacing}
\usepackage{graphicx}
% Circuitos: el trabajo de electrónica se hace hoy con micro:bit y sus sensores
% integrados (MakeCode, sin cableado), así que no se necesitan esquemáticos.
% Si algún día entran montajes con protoboard, basta descomentar esta línea:
% los diagramas .tex/.tikz declarados en `recursos:` ya se incrustan solos.
% \usepackage{circuitikz}
\usepackage[most]{tcolorbox}
\usepackage{tabularx}
\usepackage{array}
\usepackage{fancyhdr}
\usepackage{titlesec}
\usepackage[document]{ragged2e}  % bandera a la derecha en todo el cuerpo
\usepackage{enumitem}
\usepackage{microtype}
\usepackage{etoolbox}
\usepackage{needspace}
% hyperref al final del preambulo: metadatos del PDF (titulo, autor, idioma,
% familia), marcadores por estacion (\pdfbookmark) y enlaces sin recuadro.
\usepackage[hidelinks,
  pdftitle={El oficio del editor — qué hace y qué no hace},
  pdfauthor={PhD. Álvaro Cárdenas Orozco},
  pdfsubject={Tecnología e Informática · Grado 10 · Guía 1 · MILC v3},
  pdflang={es-CO},
  bookmarksopen=true,bookmarksnumbered=false]{hyperref}

% Helvetica Neue no trae la flecha U+2192, asi que cada «→» escrito literalmente
% en el YAML se compilaba a NADA, en silencio (462 flechas en 61 guias: rutas del
% tipo «paleta → tipografia → lockup» salian rotas). Mapeamos el caracter a su
% version matematica, que si tiene glifo. Asi el autor puede seguir escribiendo
% «→» comodamente en el YAML.
\usepackage{newunicodechar}
\newunicodechar{→}{\ensuremath{\rightarrow}}
% Mismo problema con los simbolos de marca y vinetas: el «✓» de «una palomita ✓»
% salia como cuadrito vacio en el PDF de todas las guias que lo usaban.
\usepackage{pifont}
\newunicodechar{✓}{\ding{51}}
\newunicodechar{✗}{\ding{55}}
\newunicodechar{✔}{\ding{52}}
\newunicodechar{•}{\textbullet}
\newunicodechar{≥}{\ensuremath{\geq}}
\newunicodechar{≤}{\ensuremath{\leq}}

\setmainfont{Helvetica Neue}
\setsansfont{Helvetica Neue}
\setlength{\parindent}{0pt}
\setlength{\parskip}{6pt}
% Sin particion de palabras: en 6.º-7.º una palabra cortada por guion al final
% de linea («Comuni-cacion») frena la lectura mucho mas que una linea corta.
\hyphenpenalty=10000
\exhyphenpenalty=10000
\pretolerance=2000
\tolerance=3000
\setlength{\emergencystretch}{3em}
\linespread{1.10}
\renewcommand{\arraystretch}{1.25}
\setlist{leftmargin=5mm,itemsep=1mm,topsep=1mm}
\newcolumntype{Y}{>{\RaggedRight\arraybackslash}X}

\definecolor{milcMagenta}{HTML}{E6007E}
\definecolor{milcVino}{HTML}{5A0038}
\definecolor{milcTurquesa}{HTML}{0093A5}
\definecolor{milcVerde}{HTML}{58A923}
\definecolor{milcMostaza}{HTML}{E5B400}
\definecolor{milcOcre}{HTML}{B86600}
\definecolor{milcNegro}{HTML}{241816}
\definecolor{milcGris}{HTML}{F7F7F4}
\definecolor{milcLinea}{HTML}{E8E2DD}
\definecolor{milcGradePrimary}{HTML}{FF2D87}
\definecolor{milcGradeDark}{HTML}{9F1B56}
\definecolor{milcGradeSoft}{HTML}{FFE0EE}
\definecolor{milcGradeAccent}{HTML}{CFE7FF}
\definecolor{milcGradeText}{HTML}{FFFFFF}
\color{milcNegro}

\pagestyle{fancy}
\fancyhf{}
% Cabecera de dos lineas: identidad de la guia arriba y, a la derecha y en
% gris, la estacion en curso (\markright desde \milcestacion). Antes de la
% primera estacion la segunda linea va vacia; el \strut mantiene la altura
% para que la linea de arriba no baile entre paginas.
\lhead{\small Tecnología e Informática · Grado 10\\[-1pt]{\footnotesize\strut}}
\rhead{\small Guía 1 · MILC v3\\[-1pt]{\footnotesize\strut\color{milcNegro!65}\rightmark}}
\cfoot{\small\thepage}
\renewcommand{\headrulewidth}{0pt}

% Sobre mostaza (#E5B400) y verde (#58A923) el texto va en milcNegro; sobre el
% resto de la paleta, en blanco. Se usa dentro de fonttitle/fontupper, donde un
% \color posterior manda sobre coltitle.
\newcommand{\milctintasobre}[1]{%
  \ifboolexpr{test{\ifstrequal{#1}{milcMostaza}} or test{\ifstrequal{#1}{milcVerde}}}%
    {\color{milcNegro}}{\color{white}}}

\newtcolorbox{titlebox}[1]{
  enhanced,colback=#1,colframe=#1,arc=7mm,boxrule=0pt,
  left=7mm,right=7mm,top=3mm,bottom=3mm,
  before skip=8pt,after skip=8pt,
  fontupper=\sffamily\bfseries\Large\color{white}
}

\newtcolorbox{softbox}[3]{
  enhanced,breakable,pad at break=3mm,
  colback=#2,colframe=#1,borderline west={4pt}{0pt}{#1},
  arc=2mm,boxrule=.4pt,left=4mm,right=4mm,top=3mm,bottom=3mm,
  before skip=6pt,after skip=8pt,title={#3},coltitle=white,
  colbacktitle=#1,fonttitle=\sffamily\bfseries\milctintasobre{#1},fontupper=\RaggedRight,
  attach boxed title to top left={xshift=3mm,yshift=-2mm},
  boxed title style={arc=2mm, boxrule=0pt}
}

\newtcolorbox{drawbox}[2]{
  enhanced,breakable,pad at break=4mm,
  colback=white,colframe=#1,arc=4mm,boxrule=1pt,
  left=5mm,right=5mm,top=4mm,bottom=4mm,before skip=8pt,after skip=8pt,
  title={#2},coltitle=white,colbacktitle=#1,fonttitle=\sffamily\bfseries\milctintasobre{#1},
  fontupper=\RaggedRight,
  attach boxed title to top left={xshift=4mm,yshift=-2mm},
  boxed title style={arc=3mm, boxrule=0pt}
}

\newtcolorbox{infoband}[1]{
  enhanced,breakable,pad at break=2mm,colback=#1!8,colframe=#1!50,arc=2mm,boxrule=.5pt,
  left=4mm,right=4mm,top=2mm,bottom=2mm,before skip=4pt,after skip=4pt,
  fontupper=\small\RaggedRight
}

% Puente narrativo entre secciones (contrato editorial Conéctate).
% Da continuidad: "Acabas de X. Ahora vas a Y porque Z."
% Visualmente: banda izquierda en color del grado, texto en itálica pequeña.
\newtcolorbox{puentebox}{
  enhanced,breakable,pad at break=2mm,colback=milcGradeSoft,colframe=milcGradeDark,
  borderline west={3pt}{0pt}{milcGradeDark},
  arc=0mm,boxrule=0pt,left=5mm,right=4mm,top=2.5mm,bottom=2.5mm,
  before skip=4pt,after skip=8pt,
  fontupper=\itshape\small\color{milcNegro}\RaggedRight
}
% La flecha va en modo matematico a proposito: Helvetica Neue NO tiene el glifo
% U+2192, asi que \textrightarrow se compilaba a NADA («Missing character: There
% is no → in font Helvetica Neue») y el puente salia sin flecha. En modo
% matematico la flecha viene de la fuente matematica y si se dibuja.
\newcommand{\puente}[1]{%
  \begin{puentebox}%
    \textcolor{milcGradeDark}{$\rightarrow$}\hspace{2mm}#1%
  \end{puentebox}%
}

\newcommand{\linead}[1]{\noindent\color{milcLinea}\rule{#1}{0.5pt}\color{milcNegro}}
\newcommand{\respuesta}[1]{\par\vspace{2mm}\foreach \n in {1,...,#1}{\noindent\color{milcLinea}\rule{\linewidth}{0.5pt}\par\vspace{5mm}}\color{milcNegro}}
\newcommand{\checkbox}{\textcolor{milcGradeDark}{\fbox{\phantom{x}}}\hspace{5pt}}

% ─── Listas aireadas para las guías socioemocionales ────────────────────
% Componer las verificaciones como «\checkbox texto\\» deja el texto que salta
% de línea debajo del cuadrito y apelmaza el bloque. Estos entornos alinean la
% continuación y airean los ítems. Los usa Territorio Interior; cualquier
% familia puede hacerlo.
\newlist{checklista}{itemize}{1}
\setlist[checklista]{
  label=\textcolor{milcGradeDark}{\fbox{\phantom{x}}},
  leftmargin=9mm, labelsep=3.2mm, itemsep=2.4mm,
  topsep=2.5mm, parsep=0pt, partopsep=0pt, before=\RaggedRight
}
\newlist{pasoslista}{enumerate}{1}
\setlist[pasoslista]{
  label=\textcolor{milcGradeDark}{\sffamily\bfseries\arabic*.},
  leftmargin=8mm, labelsep=3mm, itemsep=2.8mm,
  topsep=1mm, parsep=0pt, partopsep=0pt, before=\RaggedRight
}

\newcommand{\phasecard}[4]{
\begin{tcolorbox}[enhanced,colback=#3,colframe=#2,arc=3mm,boxrule=.5pt,height=3.05cm,valign=center,halign=center,left=1.5mm,right=1.5mm,top=2mm,bottom=2mm]
{\sffamily\bfseries\fontsize{22}{24}\selectfont\textcolor{#2}{#1}}\par
{\sffamily\bfseries\small #4}
\end{tcolorbox}
}

% ── Piezas del rediseno 2026 (legibilidad 6.º-11.º) ───────────────────
% Banda de estacion: reemplaza el titlebox macizo. Titulo a la izquierda,
% dato practico (tiempo/modalidad) a la derecha, en una sola linea.
% Nunca huerfana: pide 9 lineas libres (o salta de pagina rellenando con
% \vfill) y prohibe el salto justo despues de la banda. Nueve y no seis:
% la caja partible que sigue necesita ~80pt (titulo adosado, relleno y dos
% lineas) para arrancar en la misma pagina; con menos, tcolorbox la manda
% entera a la siguiente y la banda se queda sola. El penalty 10000 va
% en `after`, ANTES del espacio posterior: un `after skip` (aunque sea 0pt)
% es un \vskip tras la caja, y ese glue es un punto de corte legal que un
% \nopagebreak escrito despues de \end{tcolorbox} ya no cierra. Cada banda
% deja marcador PDF y actualiza la cabecera derecha.
\newcounter{milcestacion}
\newcommand{\milcestacion}[2]{%
\Needspace*{9\baselineskip}%
\stepcounter{milcestacion}%
\pdfbookmark[1]{#2}{milcestacion\themilcestacion}%
\markright{#2}%
\begin{tcolorbox}[enhanced,colback=#1,colframe=#1,arc=2mm,boxrule=0pt,
  left=5mm,right=5mm,top=2.6mm,bottom=2.6mm,before skip=11pt,
  after={\par\nopagebreak\vskip7pt}]
{\sffamily\bfseries\fontsize{15}{18}\selectfont\milctintasobre{#1}#2}
\end{tcolorbox}}

% Tarjeta de paso numerada: sustituye las tablas «Pilar / Aplicacion», que
% eran formato de consulta docente, no de estudiante. El numero va en un
% circulo sobre el borde para que la secuencia se lea de un vistazo.
\newcommand{\milcpaso}[3]{%
\Needspace{4\baselineskip}%
\begin{tcolorbox}[enhanced,colback=white,colframe=milcLinea,arc=1.5mm,boxrule=.9pt,
  left=14mm,right=4mm,top=3mm,bottom=3mm,before skip=4pt,after skip=4pt,
  fontupper=\RaggedRight,
  overlay={\node[anchor=north west,circle,fill=#2,text=white,inner sep=0pt,
    minimum size=7.4mm,font=\sffamily\bfseries\small]
    at ([xshift=3.6mm,yshift=-3.2mm]frame.north west) {#1};}]
#3
\end{tcolorbox}}

% Casilla marcable de tamano util con lapiz (la anterior era \fbox{\phantom{x}},
% de alto variable segun el texto que la seguia).
\newcommand{\milcmarca}{\framebox[3.8mm]{\rule{0pt}{3.4mm}}}

% Fila marcable de lista suelta (checks de la estacion 1).
\newcommand{\milccheckrow}[1]{%
\par\vspace{1.8mm}\noindent
\makebox[6.4mm][l]{\raisebox{-.55mm}{\textcolor{milcNegro}{\milcmarca}}}%
\parbox[t]{\dimexpr\linewidth-6.4mm\relax}{\RaggedRight #1}\par}

% Celda marcable dentro de la rubrica (tabla de autoevaluacion). Misma
% sangria francesa, pero pensada para vivir dentro de una columna X.
\newcommand{\milcopcion}[1]{%
\makebox[5.2mm][l]{\raisebox{-.55mm}{\textcolor{milcNegro}{\milcmarca}}}%
\parbox[t]{\dimexpr\linewidth-5.2mm\relax}{\RaggedRight #1}}

% ── Recursos visuales de la guía ──────────────────────────────────────
% Declarados en el bloque `recursos:` del YAML y emitidos por el builder.
% \guiaFigura   → imagen rasterizada o PDF (foto, carta celeste, montaje).
% \guiaDiagrama → fuente TikZ/circuitikz (.tex) incrustada como vector:
%                 es la vía para circuitos y esquemáticos editables.
% [#1] = texto alternativo (alt) · #2 = ruta relativa al .tex · #3 = pie
% El alt llega al PDF etiquetado (/Alt) y lo lee el lector de pantalla.
\newcommand{\guiaFigura}[3][]{%
\begin{center}
\includegraphics[width=0.88\linewidth,height=0.38\textheight,keepaspectratio,alt={#1}]{#2}\\[1.5mm]
{\footnotesize\itshape\textcolor{milcNegro!75}{#3}}
\end{center}
}

\newcommand{\guiaDiagrama}[2]{%
\begin{center}
\input{#1}\\[1.5mm]
{\footnotesize\itshape\textcolor{milcNegro!75}{#2}}
\end{center}
}

\begin{document}
\thispagestyle{empty}

% ── Portada ───────────────────────────────────────────────────────────
% Antes: un rectangulo de color a pagina completa. Se veia bien en pantalla,
% pero en la fotocopiadora del colegio se comia el toner y salia gris sucio.
% Ahora la banda ocupa el tercio superior y el resto va en blanco.
\begin{tikzpicture}[remember picture,overlay]
  \path[fill=milcGradePrimary]
    (current page.north west) rectangle ([yshift=-10.6cm]current page.north east);

  \node[anchor=north west,text=milcGradeText] at ([xshift=2.0cm,yshift=-1.55cm]current page.north west)
    {\sffamily\bfseries\fontsize{12}{14}\selectfont GRADO};
  \node[anchor=north west,text=milcGradeText,opacity=.85] at ([xshift=2.0cm,yshift=-2.35cm]current page.north west)
    {\sffamily\bfseries\fontsize{8.5}{10}\selectfont SERIE GUÍAS · TECNOLOGÍA E INFORMÁTICA};
  \node[anchor=north east,text=milcGradeText,opacity=.85,align=right] at ([xshift=-2.0cm,yshift=-1.55cm]current page.north east)
    {\sffamily\bfseries\fontsize{9}{12}\selectfont I.E. Sor María Juliana\\Cartago, Valle del Cauca};

  \node[anchor=west,text=milcGradeText] (gradonum) at ([xshift=1.85cm,yshift=-7.1cm]current page.north west)
    {\sffamily\bfseries\fontsize{112}{112}\selectfont 10};
  % El circulito de grado (°) solo aplica a las guias de grado («11°»); en
  % Bebras y Semillero el numero grande es un momento/modulo, no un grado.
  % Se posiciona relativo al nodo `gradonum`, no en coordenadas absolutas.
  \draw[milcGradeText,line width=1.6pt]
    ([xshift=2.0mm,yshift=-4.5mm]gradonum.north east) circle (.15cm);

  \node[anchor=west,text=milcGradeText,text width=11.4cm,align=left] at ([xshift=1.5cm,yshift=0.5cm]gradonum.east)
    {
     {\sffamily\bfseries\fontsize{9}{11}\selectfont\MakeUppercase{Período 1 · Escritura abierta con IA}}\\[3mm]
     {\sffamily\bfseries\fontsize{21}{25}\selectfont\setlength{\baselineskip}{27pt}El oficio del editor ---\\[4pt]qué hace y\\[4pt]qué no hace}\\[3.5mm]
     {\sffamily\bfseries\fontsize{15}{18}\selectfont Guía 1-10-TIC}
    };
\end{tikzpicture}

\vspace*{9.4cm}
\vfill

{\sffamily\bfseries\fontsize{8}{10}\selectfont\textcolor{milcGradeDark}{LO QUE TE LLEVAS HOY}}

\begin{tcolorbox}[enhanced,colback=white,colframe=milcNegro,arc=2mm,boxrule=1.6pt,
  left=5mm,right=5mm,top=4mm,bottom=4mm,before skip=3pt,after skip=12pt,
  fontupper=\RaggedRight\fontsize{11}{15.5}\selectfont]
Definición personal del rol de editor en 1 página, las 3 decisiones críticas declaradas, el análisis de 2 libros con lente editorial y un cierre sobre por qué editar no es escribir
\end{tcolorbox}

\begin{tcolorbox}[enhanced,colback=milcGris,colframe=milcGris,arc=2mm,boxrule=0pt,
  left=5mm,right=5mm,top=3.5mm,bottom=3.5mm,after skip=12pt]
{\sffamily\bfseries\fontsize{8}{10}\selectfont\textcolor{milcNegro!55}{ESTA GUÍA ES DE}}\\[3mm]
\begin{tabularx}{\linewidth}{X p{3.2cm} p{3.2cm}}
\linead{\linewidth} & \linead{3.0cm} & \linead{3.0cm}\\[-1mm]
{\fontsize{8.5}{10}\selectfont\textcolor{milcNegro!55}{Nombre completo}} &
{\fontsize{8.5}{10}\selectfont\textcolor{milcNegro!55}{Grupo}} &
{\fontsize{8.5}{10}\selectfont\textcolor{milcNegro!55}{Fecha}}\\
\end{tabularx}
\end{tcolorbox}

\vfill

\noindent\textcolor{milcLinea}{\rule{\linewidth}{1pt}}\\[2mm]
{\sffamily\bfseries\fontsize{9.5}{12}\selectfont Autor: Dr. Álvaro Cárdenas Orozco}\\[1mm]
{\sffamily\fontsize{8.5}{11}\selectfont\textcolor{milcNegro!55}{Metodología MILC · apertura ancestral · pensamiento computacional · triángulo Dussel–estoicismo–Floridi}}

\newpage

\section*{Apertura: El oficio del editor --- qué hace y qué no hace}

\begin{softbox}{milcGradeDark}{milcGradeSoft}{Saber ancestral}
En Cartago la reputación de un taller de bordado no la construye lo que se dice de él, sino lo que sale de sus manos. Los ornamentos que vistió el papa durante su visita a Colombia se bordaron allí (Chica García, 2019). Una prenda pasa por unas ocho artesanas y más de ocho horas de trabajo antes de llevar el nombre de un taller. En un oficio así el nombre no se defiende hablando: se defiende con la puntada, y una sola pieza mal hecha pesa más que cualquier rumor. La \textbf{cara de exclusión} está en ese mismo nombre. El taller lo sostienen ocho artesanas por prenda que rara vez aparecen en él, porque la reputación se acumula en el taller y no en quien borda. Firmar es quedar respondiendo por lo que otras manos hicieron. Y también es decidir a quién se nombra.\par\smallskip{\footnotesize\textcolor{milcNegro!70}{\textbf{Fuente.} Chica García, A. (2019, 17 de agosto). \emph{Bordados de Cartago: la herencia española que apropiaron las mujeres vallunas}. Radio Nacional de Colombia.}}
\end{softbox}

\begin{softbox}{milcTurquesa}{milcGris}{Saber contemporáneo}
El oficio editorial existe desde que hubo libros, mucho antes de internet y de la IA. Un \textbf{editor} es quien toma las decisiones sobre una obra: qué entra y qué sale, cómo se ordena, para quién es, qué tono tiene, qué se acepta y qué se devuelve a corregir. En el oficio se dice corto: el editor no es el autor, pero firma como responsable de la pieza. En este periodo tú vas a ser el editor de tu propio libro. La IA generativa será tu escritor asistente y producirá primeros borradores. Tú decides, ordenas, refinas y firmas. Si pierdes esa distinción terminas con un libro que la IA escribió y tú presentaste. Si la sostienes, terminas con un libro tuyo, hecho con ayuda de IA.
\end{softbox}

\begin{softbox}{milcMostaza}{milcGris}{Pregunta-puente}
\textit{¿Qué está aceptando un taller cuando pone su nombre en una prenda que pasó por ocho manos? ¿Y qué acepta quien firma un libro escrito con ayuda de una IA?}
\end{softbox}

\begin{softbox}{milcVerde}{milcGris}{Saber y saber hacer}
Al terminar podrás: (1) \textbf{identificar} las tres decisiones que cualquier editor toma sobre una obra --- tema, audiencia y estructura --- en libros que ya conoces; (2) \textbf{analizar} qué decisión editorial sostiene principalmente a cada uno de esos libros; (3) \textbf{evaluar} qué responsabilidad estás aceptando al firmar como editor de tu propio libro.
\end{softbox}



\small
\begin{tabularx}{\textwidth}{>{\bfseries}p{3.0cm}Y}
\rowcolor{milcGradePrimary!12}Autor & Dr. Álvaro Cárdenas Orozco\\
DBA & Producción editorial mediada por IA con criterio ético y técnico (MEN, T\&I)\\
Referentes & Ley 23 de 1982 · Floridi (ética de la IA generativa) · Dussel (autoría con dignidad)\\
Producto final & Definición personal del rol de editor en 1 página, las 3 decisiones críticas declaradas, el análisis de 2 libros con lente editorial y un cierre sobre por qué editar no es escribir\\
\end{tabularx}
\normalsize

\puente{Este periodo te lleva a producir un libro como editor responsable, con la IA generativa de escritor asistente. Las nueve sesiones que siguen se sostienen sobre la decisión que tomas hoy: aceptar el rol de editor.}

\milcestacion{milcGradeDark}{Ruta MILC: así vamos a aprender}
\begin{tabularx}{\textwidth}{>{\centering\arraybackslash}X>{\centering\arraybackslash}X>{\centering\arraybackslash}X>{\centering\arraybackslash}X}
\phasecard{1}{milcVerde}{milcGris}{Escucha\\[-1mm]\small Observo, escucho y pregunto.} &
\phasecard{2}{milcTurquesa}{milcGris}{Entender\\[-1mm]\small Sistematizo y ordeno.} &
\phasecard{3}{milcMagenta}{milcGris}{Hacer\\[-1mm]\small Construyo y aplico.} &
\phasecard{4}{milcVino}{milcGris}{Cierre\\[-1mm]\small Evalúo y reflexiono.}
\end{tabularx}


\puente{Antes de teorizar, vas a mirar dos libros que conoces de verdad. De cada uno vas a anotar quién lo escribió, para quién parece hecho, cómo está organizado y por qué te funcionó a ti.}

\milcestacion{milcVerde}{Estación 1 de 3 · Escucha}

\begin{softbox}{milcVerde}{milcGris}{Escena para escuchar}
\textbf{Actividad 1 · IDENTIFICA --- Dos libros bajo la lupa} (15 min · individual). (1) Anota dos libros que conozcas de verdad: un manga, una novela, un manual, un libro de divulgación o un cómic. (2) Para cada uno responde sin buscar nada: quién es el autor o autora. (3) Para quién parece hecho el libro, y en qué lo notas. (4) Cómo está organizado por dentro: capítulos, secciones, ilustraciones. (5) Qué hace que a ti te funcione, en una frase.
\end{softbox}

{\sffamily\bfseries\fontsize{8}{10}\selectfont\textcolor{milcNegro!55}{MARCA CUANDO LO HAYAS HECHO}}
\milccheckrow{Elegí dos libros que conozco bien, no dos que me suenan.}
\milccheckrow{Respondí las cuatro preguntas sobre cada uno sin buscar información.}
\milccheckrow{Escribí en una frase qué hace que cada libro me funcione a mí.}

\begin{infoband}{milcVerde}


\textbf{Lo que va al cuaderno (Actividad 1 · IDENTIFICA).} \textbf{Título:} «Actividad 1 --- Dos libros bajo la lupa». \textbf{Formato:} una ficha por libro con autoría, audiencia, organización y la frase de por qué te funciona. \textbf{Extensión:} un tercio de página. \textbf{Sabes que terminaste cuando} ves que cada libro es resultado de decisiones y no de casualidades.
\end{infoband}

\puente{Ya miraste dos libros. Ahora vas a entender la \textbf{anatomía del oficio}: las tres decisiones críticas, los errores típicos del editor novato con IA y la diferencia entre escribir y editar.}

\milcestacion{milcTurquesa}{Estación 2 de 3 · Entender}

\begin{softbox}{milcTurquesa}{milcGris}{Lo que vas a entender}
\textbf{Actividad 2 · ANALIZA --- Las tres decisiones y los cinco errores} (20 min · parejas). (1) Con tu pareja, escriban las tres decisiones críticas con un ejemplo propio de cada una. (2) Escriban los cinco errores del editor novato con IA. (3) Cada uno marca el error que más teme cometer y explica por qué. (4) Escriban en dos líneas la diferencia entre escribir y editar, con sus propias palabras.
\end{softbox}

\milcpaso{1}{milcTurquesa}{Las tres decisiones más críticas se descomponen así. El \textbf{tema}: de qué trata el libro y qué pregunta responde, en una o dos frases. Si no cabe en dos frases, el tema todavía no está claro. La \textbf{audiencia}: para quién es, qué edad tiene y cuánto conocimiento previo se le supone. La \textbf{estructura}: cómo se ordena el contenido, cuántos capítulos, si es narrativo o temático. Las tres se toman antes de escribir la primera palabra.}
\milcpaso{2}{milcTurquesa}{La diferencia entre escribir y editar cabe en una imagen: quien escribe llena la página en blanco, y quien edita decide qué se queda en la página llena. Producir texto es justo lo que la IA hace con facilidad. Decidir qué texto merece quedarse sigue siendo trabajo humano, porque exige criterio cultural, ético y estético. La phronesis del editor consiste en no avergonzarse de pedirle a la IA, y en no delegarle la decisión.}
\milcpaso{3}{milcTurquesa}{Los cinco errores del editor novato con IA. Aceptar el primer borrador, que casi nunca es el definitivo. Delegarle a la IA las tres decisiones críticas. No leer lo que se firma. Ocultar que se usó IA. Y renunciar a la voz propia, con lo que el libro termina sonando a chatbot y no a ti.}
\milcpaso{4}{milcTurquesa}{Algoritmo del oficio editorial con IA: (1) toma las tres decisiones críticas tú; (2) genera borradores con la IA; (3) lee cada borrador como crítico y no como receptor; (4) reescribe lo que no suene a ti; (5) firma solo lo que puedes defender; (6) declara qué aportó la IA.}

\begin{drawbox}{milcTurquesa}{Anatomía del oficio editorial}
\textbf{Tema:} de qué trata, en una o dos frases. \\[1mm] \textbf{Audiencia:} para quién, con qué nivel previo. \\[1mm] \textbf{Estructura:} cómo se ordena el contenido. \\[1mm] \textbf{Regla de oro:} no escribes todo, pero firmas todo. \\[1mm] \textbf{Con IA:} la IA propone, tú decides, tú firmas y tú declaras.
\end{drawbox}

\begin{softbox}{milcMostaza}{milcGris}{Errores comunes y su corrección}
(1) Aceptar el primer borrador sin editarlo. (2) Delegarle a la IA las tres decisiones críticas. (3) No leer lo que se firma. (4) Ocultar el uso de IA en el producto final. (5) Renunciar a la voz propia y dejar que el libro suene a chatbot.
\end{softbox}

\begin{infoband}{milcTurquesa}


\textbf{Lo que va al cuaderno (Actividad 2 · ANALIZA).} \textbf{Título:} «Actividad 2 --- Anatomía editorial». \textbf{Formato:} las tres decisiones con tu propio ejemplo, los cinco errores con el que más temes marcado, y la diferencia entre escribir y editar en dos líneas. \textbf{Extensión:} media página. \textbf{Sabes que terminaste cuando} podrías explicarle a un compañero qué hace un editor de libro.
\end{infoband}

\puente{Tienes el método. Toca aplicarlo. Escribes tu definición de editor, declaras tus tres decisiones, analizas los dos libros con esa lente y cierras con la diferencia entre escribir y editar.}

\milcestacion{milcMagenta}{Estación 3 de 3 · Hacer}

\begin{softbox}{milcMagenta}{milcGris}{Cómo vas a construir tu producto}
\textbf{Actividad 3 · EVALÚA --- Mi definición de editor} (30 min · individual). (1) Escribe a mano, sin pedirle nada a la IA, qué te comprometes a hacer como editor durante el periodo. (2) Declara tus tres decisiones preliminares: tema, audiencia y estructura. (3) Analiza los dos libros de la actividad 1 y di qué decisión editorial sostiene principalmente a cada uno. (4) Cierra con cinco líneas sobre por qué editar no es escribir. (5) Lee tu definición en voz alta y pregúntate si suena a ti o a respuesta de chatbot. (6) Fírmala con tu nombre.
\end{softbox}

\milcpaso{1}{milcMagenta}{Tu entrega se construye en cuatro piezas: la definición personal de una página, las tres decisiones críticas declaradas, el análisis de los dos libros y el párrafo de cierre. Las cuatro son tuyas, sin excepción.}
\milcpaso{2}{milcMagenta}{Sigue el patrón «voz propia desde el principio». La definición se escribe a mano, sin IA. Es una declaración tuya sobre el oficio. Si desde la primera sesión la IA escribe por ti, vas a llegar a la última sin haber asumido el rol.}
\milcpaso{3}{milcMagenta}{Las tres decisiones pueden ser preliminares. No tienes que cerrar hoy el tema del libro. Pero declararlas obliga a pensar antes de actuar, y eso es distinto de abrir un chat y preguntar qué libro hacer.}
\milcpaso{4}{milcMagenta}{Algoritmo: (1) escribe la definición; (2) declara las tres decisiones; (3) analiza los dos libros con la lente editorial; (4) cierra con la diferencia entre escribir y editar; (5) léela en voz alta; (6) fírmala.}

\begin{drawbox}{milcMagenta}{Checklist antes de entregar}
\checkbox Definición personal de una página, escrita sin IA. \\[1mm] \checkbox Las tres decisiones críticas declaradas, aunque sean preliminares. \\[1mm] \checkbox Los dos libros analizados con lente editorial. \\[1mm] \checkbox Párrafo de cierre sobre escribir y editar. \\[1mm] \checkbox Leída en voz alta y reconocible como tu voz. \\[1mm] \checkbox Firmada con tu nombre.
\end{drawbox}

\begin{softbox}{milcMagenta}{milcGris}{Plantilla de trabajo}
\textbf{Definición personal.} \textit{Qué me comprometo a hacer como editor.} \\[1mm] \textbf{Decisión 1 (tema).} \textit{De qué tratará mi libro.} \\[1mm] \textbf{Decisión 2 (audiencia).} \textit{Para quién es.} \\[1mm] \textbf{Decisión 3 (estructura).} \textit{Cómo se organiza.} \\[1mm] \textbf{Análisis de los dos libros.} \textit{Qué decisión editorial sostiene a cada uno.} \\[1mm] \textbf{Cierre.} \textit{Por qué editar no es escribir.}
\end{softbox}

\begin{infoband}{milcMagenta}


\textbf{Lo que va al cuaderno (Actividad 3 · EVALÚA).} \textbf{Título:} «Actividad 3 --- Mi definición de editor». \textbf{Formato:} la definición, las tres decisiones, los dos análisis y el cierre, firmados con tu nombre. \textbf{Extensión:} una página. \textbf{Sabes que terminaste cuando} podrías firmar ese texto con tu nombre sin incomodidad.
\end{infoband}

\puente{Lo que sigue resume \emph{qué} entregas hoy: la definición, las tres decisiones, los dos análisis y el cierre. El criterio principal es que al terminar sepas qué responsabilidad estás aceptando para las nueve sesiones que vienen.}

\milcestacion{milcMostaza}{Producto de la sesión}
\textbf{Definición personal de editor, 3 decisiones críticas, análisis de 2 libros y cierre}.
\begin{softbox}{milcMostaza}{milcGris}{Criterios de calidad}
(1) Definición personal de una página escrita sin asistencia de IA. (2) Las tres decisiones críticas declaradas, aunque sean preliminares. (3) Los dos libros analizados con lente editorial, no resumidos. (4) Párrafo de cierre sobre la diferencia entre escribir y editar. (5) Voz propia reconocible al leerla en voz alta. (6) Entrega firmada con tu nombre.
\end{softbox}

\puente{Antes de cerrar, mira el oficio del editor desde las cinco dimensiones humanas. Aceptar el rol con honestidad es respeto por el oficio. Declararse «solo usuario de IA» es renunciar a la responsabilidad antes de empezar.}

\milcestacion{milcVino}{Evaluación liberadora · cinco dimensiones}

\begin{softbox}{milcVino}{milcGris}{Sentido de la evaluación}
Evaluar no es solo revisar una nota. Es reconocer qué aprendí, qué cambió en mí, qué emociones manejé, cómo me relaciono con la ciudad y la comunidad, y qué le dejaré a quien viene detrás.
\end{softbox}

\textbf{Mi reflexión liberadora en cinco dimensiones.} Completa cada frase iniciada:

\small
\begin{tabularx}{\textwidth}{>{\bfseries\RaggedRight\arraybackslash}p{3.4cm}Y>{\RaggedRight\arraybackslash}p{4.6cm}}
\rowcolor{milcVino}\color{white}Dimensión & \color{white}\bfseries Mi respuesta (frame starter) & \color{white}\bfseries Algo que descubrí\\
1. Personal & Lo que más cambió en mí fue \linead{6.5cm} & \linead{4.2cm}\\
2. Emocional & La emoción más fuerte fue \linead{2.5cm} y la regulé \linead{3cm} & \linead{4.2cm}\\
3. Ciudadana & En democracia esto importa porque \linead{6cm} & \linead{4.2cm}\\
4. Local (Cartago) & En mi barrio sería útil para \linead{6.5cm} & \linead{4.2cm}\\
5. Intergeneracional & A mi abuela/abuelo le diría \linead{6.5cm} & \linead{4.2cm}\\
\end{tabularx}
\normalsize

\begin{infoband}{milcVino}
\textbf{Para la próxima sustentación.} Vuelve a esta tabla en seis meses. Verás que las respuestas cambian — no porque lo aprendido cambie, sino porque tú cambias. La evaluación liberadora es una conversación contigo a través del tiempo.
\end{infoband}


\puente{Y para terminar, tres voces sobre firmar lo que otras manos hicieron. Dussel sobre el trabajo que hay detrás de cada objeto, Séneca sobre que la vida no desdiga de la palabra, Floridi sobre artefactos que ya no solo obedecen instrucciones.}

\milcestacion{milcGradeDark}{Triángulo de pensamiento · cierre filosófico}

\begin{softbox}{milcVino}{milcGris}{Dussel · lente del nosotros}
\textit{«Los artefactos u objetos… que nos rodean en el mundo cotidiano… son entonces productos del trabajo humano… Si los entes son instrumentos producidos significa que fueron objeto de un cierto trabajo.»}

{\footnotesize\itshape\color{milcNegro!70}--- Enrique Dussel · Filosofía de la liberación (1977), §4.3.4.3}

Detrás de una prenda bordada hay ocho artesanas y más de ocho horas. Detrás de un libro hay lecturas, decisiones y descartes. Firmar como editor no es afirmar que lo escribiste todo: es responder por el trabajo que hay debajo, incluido el que no hiciste tú.

\textbf{De lo que voy a firmar este periodo, ¿qué parte va a ser trabajo mío y qué parte de otros?}
\end{softbox}
\linead{\linewidth}\\[1mm]
\linead{\linewidth}

\begin{softbox}{milcMostaza}{milcGris}{Estoicismo · Séneca · Cartas a Lucilio, 20 (c. 64 d.C.) · cuidado interior}
\textit{«La filosofía enseña a hacer, no a decir; y exige que cada uno viva según su propia ley, que la vida no desdiga de la palabra. (trad. propia)»}

{\footnotesize\itshape\color{milcNegro!70}--- Séneca · Cartas a Lucilio, 20 (c. 64 d.C.)}

El taller de Cartago no defiende su nombre hablando: lo defiende con la puntada. Tu definición de editor no vale por lo bien redactada que quede hoy, sino por si dentro de nueve sesiones el libro se le parece. Lo que escribas hoy va a quedar como medida de lo que hagas después.

\textbf{Lo que escribí hoy que me comprometo a hacer, ¿lo estoy haciendo en la sesión 5?}
\end{softbox}
\linead{\linewidth}\\[1mm]
\linead{\linewidth}

\begin{softbox}{milcTurquesa}{milcGris}{Floridi · lente de la infoesfera}
\textit{«En el mundo onlife, los artefactos han dejado de ser meras máquinas que simplemente operan según instrucciones humanas: pueden cambiar de estado de manera autónoma. (trad. propia)»}

{\footnotesize\itshape\color{milcNegro!70}--- The Onlife Initiative (ed. Luciano Floridi) · The Onlife Manifesto (2015), § 3.2}

Una herramienta que propone texto por su cuenta no es un martillo. Por eso la pregunta de quién firma se volvió difícil justo ahora. La respuesta del oficio sigue siendo la misma que en el taller: firma quien decide, y quien decide responde.

\textbf{Cuando la IA me proponga un capítulo entero, ¿voy a decidir sobre él o voy a aceptarlo?}
\end{softbox}
\linead{\linewidth}\\[1mm]
\linead{\linewidth}

\begin{infoband}{milcNegro}
\textbf{Nota para el docente.} Las tres son citas textuales y llevan su referencia debajo. Puedes remitir a la obra en clase.
\end{infoband}

\milcestacion{milcVino}{Mi autoevaluación · marca una casilla por fila}

{\small
\setlength{\extrarowheight}{2.2pt}
\begin{tabularx}{\textwidth}{>{\RaggedRight\arraybackslash\bfseries}p{3.0cm}YYY}
\rowcolor{milcVino}\color{white}\bfseries Criterio & \color{white}\bfseries Lo logré & \color{white}\bfseries En proceso & \color{white}\bfseries Necesito apoyo\\[.6mm]
Comprensión del tema & \milcopcion{Explico las ideas con mis palabras.} & \milcopcion{Reconozco las ideas, falta precisión.} & \milcopcion{Necesito repasar lo básico.}\\[.8mm]
\rowcolor{milcGris}Pensamiento computacional & \milcopcion{Usé los pilares al armar mi producto.} & \milcopcion{Usé algunos pilares.} & \milcopcion{Necesito releer la estación 2.}\\[.8mm]
Aplicación y producto & \milcopcion{Mi producto cumple los criterios.} & \milcopcion{Está armado, con ajustes pendientes.} & \milcopcion{Aún está en construcción.}\\[.8mm]
\rowcolor{milcGris}Reflexión 5 dimensiones & \milcopcion{Respondí cada una con honestidad.} & \milcopcion{Respondí algunas con detalle.} & \milcopcion{Mis respuestas son muy generales.}\\[.8mm]
Triángulo de pensamiento & \milcopcion{Conecté las 3 voces con mi trabajo.} & \milcopcion{Conecté al menos una.} & \milcopcion{No las relaciono con mi proceso.}\\[.8mm]
\end{tabularx}}

\vspace{3mm}
\textbf{Hoy entendí que:}\\[1mm]
\linead{\linewidth}\\[1mm]
\linead{\linewidth}

\textbf{Mi compromiso para la próxima sesión es:}\\[1mm]
\linead{\linewidth}\\[1mm]
\linead{\linewidth}

\begin{softbox}{milcMostaza}{milcGris}{Cita de despedida · Marco Aurelio}
\textit{``El obstáculo en el camino se vuelve el camino. Lo que estorba al camino se vuelve camino.''}\\[1mm]
Cada error de hoy, bien observado, se convierte en pieza del camino que pisarás mañana.
\end{softbox}

\small
\begin{tabularx}{\textwidth}{>{\bfseries}p{3.6cm}Y}
\rowcolor{milcGradeDark!12}Nombre completo: & \linead{10cm}\\
Fecha: & \linead{4cm} \quad Curso/grupo: \linead{4cm}\\
Mi firma: & \linead{9cm}\\
\end{tabularx}
\normalsize

\begin{infoband}{milcGradeDark}
\textbf{Para la próxima sesión.} Llega con tu definición de editor, las tres decisiones preliminares y el análisis de los dos libros. En la sesión 2 vas a concebir tu libro: género, audiencia precisa y escaleta. La decisión de hoy es la base de todo lo que sigue.
\end{infoband}

\end{document}
